% ==========================================================
% STATUS, COSTS, PATH, DISCUSSION
% ==========================================================

\divider{Where \vtwo{} actually is}{what exists, what does not, what I chose not to build}

% ----------------------------------------------------------
\begin{frame}{Implemented today vs.\ not yet}
    \begin{columns}[T]
        \begin{column}{0.5\textwidth}
            \begin{exampleblock}{Runs today (simulator only)}
                \footnotesize
                \begin{itemize}
                    \item Controller: authority lock, scheduler, durable attempts, dispatch blocks, recovery
                    \item SQLite store: migrations, revisions, events, lease, idempotency records;
                          backup/restore/migrate via \code{qserver-v2-admin}
                    \item OIDC JWT authentication, three scopes, HTTPS + SSE, OpenAPI
                    \item Worker protocol: heartbeat, correlation, safe stop, orphan stop, lock and fencing
                    \item Worker SDK (\code{registry.register}), profile-mode loader, worker revision
                    \item One operation: \code{simulated-count} v1 on \code{ophyd.sim}; systemd unit, example configs
                    \item 81 behavioural tests in \code{tests/v2/} (contracts, storage, scheduler, recovery,
                          web, worker runtime, admin), including the fault matrix
                \end{itemize}
            \end{exampleblock}
        \end{column}
        \begin{column}{0.48\textwidth}
            \begin{alertblock}{Not yet}
                \footnotesize
                \begin{itemize}
                    \item Any real beamline adapter, test IOC run, or hardware approval
                    \item Operator UI, client SDK, console-output streaming
                    \item Pause/resume/stop/abort/halt as a user-facing surface; today only
                          \emph{safe stop at next checkpoint} and \emph{stop after current}
                    \item Loop mode, instructions, batch submission
                    \item Interactive kernel, script upload, function execution (deliberately; next slide)
                    \item Compatibility with \code{bluesky-queueserver-api}, \code{bluesky-widgets},
                          \code{bluesky-httpserver}
                    \item HA, PostgreSQL, SQLite over NFS, worker reattachment, automatic retry
                \end{itemize}
            \end{alertblock}
            \vspace{0.1cm}
            {\scriptsize\color{qsgray} Size, for orientation: \code{v2/} is about 7.3k lines;
            \code{manager/} about 20k; \code{bluesky-httpserver} and \code{bluesky-queueserver-api}
            roughly 6k each excluding docstrings. Smaller because it does less, on purpose.\par}
        \end{column}
    \end{columns}
\end{frame}

% ----------------------------------------------------------
\begin{frame}{Things I deliberately did not build}
    \small
    \renewcommand{\arraystretch}{1.15}
    \begin{tabular}{@{}L{4.2cm}L{9.0cm}@{}}
        \toprule
        \textbf{Not built} & \textbf{Why} \\
        \midrule
        A generic \code{execute\_plan(name, args, kwargs)} & It recreates the namespace-as-API surface. Everything else in \vtwo{} is downstream of refusing it. \\
        0.x protocol or queue compatibility & Compatibility keeps the dynamic execution, transport and storage contracts that are the reason for the rewrite. Coexistence instead: separate state, endpoints and authority. \\
        Controller-managed environment updates (install, git pull, env open/close) & Deployment changes belong to deployment tooling with their own audit; the worker revision records what ran. \\
        Watchdog and worker reattachment & Independent lifetimes are where the two views of the truth diverge. \vtwo{} fails closed and asks a human. \\
        Redis & SQLite gives the transaction, revision and event boundary one controller needs, on a local disk, with no second service. \\
        Bluesky inside the controller & It couples the public authority to beamline dependencies and plan failures. \\
        Active-active or HA controllers & Leadership and fencing across hosts is not justified for one instrument; one authority is the point. \\
        Bluesky document storage & Documents stay in the RunEngine subscriptions; the controller keeps run UIDs only. \\
        \bottomrule
    \end{tabular}
\end{frame}

% ----------------------------------------------------------
\begin{frame}{Costs and risks, honestly}
    \begin{columns}[T]
        \begin{column}{0.5\textwidth}
            \textbf{What beamlines pay}
            \small
            \begin{itemize}
                \item \textbf{Adapters are work.} Every exposed workflow needs a request model, a device map
                      and a reviewed registration. Boilerplate risk is real; offline generation from
                      \code{existing\_plans\_and\_devices.yaml} and annotations can draft them, not commit them.
                \item \textbf{Less freedom at the API.} No ad-hoc plan, no script upload, no function call
                      from a GUI. Interactive Bluesky stays available for that; it just does not share
                      authority with \vtwo{}.
                \item \textbf{An identity provider is required.} Today's single-user API key is not an option.
                \item \textbf{Local disk only.} The controller host owns the state; moving it is a backup/restore.
                \item \textbf{Operators see new words}: take control, release control, stop after current,
                      recovery required. The UI must hide lease IDs and revisions.
            \end{itemize}
        \end{column}
        \begin{column}{0.48\textwidth}
            \textbf{What the project pays}
            \small
            \begin{itemize}
                \item \textbf{Two codebases during transition.} 0.x needs a stated maintenance policy
                      while \vtwo{} matures; I do not think both should be developed indefinitely.
                \item \textbf{The ecosystem does not follow for free.} \code{bluesky-widgets},
                      \code{queue-monitor}, agents built on \code{bluesky-queueserver-api} need a \vtwo{} client.
                \item \textbf{My failure matrix may be wrong.} It is tested against a fault-injecting
                      worker and the simulator, not against years of beamline incidents. That knowledge is
                      in this room.
                \item \textbf{Single author so far.} Review by people who know why 0.x made each choice
                      is the fastest way to find what I missed.
            \end{itemize}
        \end{column}
    \end{columns}
\end{frame}

% ----------------------------------------------------------
\begin{frame}{A path, not a cliff}
    \begin{columns}[T]
        \begin{column}{0.5\textwidth}
            \small
            \begin{enumerate}
                \item \textbf{Agree on the properties first.} If the eight properties are right, they constrain
                      any implementation, including a hypothetical 0.x successor that is not \vtwo{}.
                \item \textbf{Review \vtwo{} against them.} Code, tests and the failure matrix are in-tree;
                      the simulator demo runs on a laptop.
                \item \textbf{One workflow, one beamline, simulation or test IOC.} Load the real profile in
                      profile mode, write one adapter, run the failure matrix against it.
                \item \textbf{Coexist.} 0.x keeps serving what has not moved. No shared queue, endpoint,
                      worker or state; one authority per instrument at a time.
                \item \textbf{Decide the 0.x policy together}: security fixes only, or feature freeze with a date.
                \item \textbf{Only then} talk about hardware, a UI, and a client library.
            \end{enumerate}
        \end{column}
        \begin{column}{0.48\textwidth}
            \begin{block}{What would change my mind}
                \small
                \begin{itemize}
                    \item A property on the list that you can show is wrong for beamline operations.
                    \item A workflow that cannot be expressed as a catalogued operation without losing
                          something scientifically essential.
                    \item Evidence that 0.x can be brought to these properties incrementally without breaking
                          its public API anyway.
                \end{itemize}
            \end{block}
            \vspace{0.2cm}
            {\small\color{qsgray} My preference is clear: converge on one implementation once \vtwo{} has
            earned it, rather than maintain two. I would rather be argued out of that in this room than
            discover the reason at a beamline.\par}
        \end{column}
    \end{columns}
\end{frame}

% ----------------------------------------------------------
\begin{frame}{Questions I would like this group to answer}
    \normalsize
    \begin{enumerate}
        \setlength{\itemsep}{0.25cm}
        \item Which of the eight properties do you disagree with, and why?
        \item Which 0.x capability did I drop that you would fight to keep
              (interactive kernel, script upload, function execution, loop mode, \ldots)?
        \item Is ``catalog, not namespace'' workable for the beamlines you support?
              Which workflow would you pilot first?
        \item Where is the failure matrix wrong? What incident would it have handled badly?
        \item We intend to expose this service on the internet behind facility SSO.
              Which of the 0.x assumptions on the security slide would you be comfortable keeping there?
        \item If \vtwo{} proceeds, what is the right maintenance policy for 0.x?
    \end{enumerate}
    \vspace{0.4cm}
    \begin{center}
        {\normalsize\color{qsgray} Code: \code{src/bluesky\_queueserver/v2/} \quad
        Tests: \code{src/bluesky\_queueserver/tests/v2/} \quad
        Docs: \code{docs/source/queueserver\_v2.rst} \quad
        Demo: \code{demo/}\\[0.2cm]
        \footnotesize Backup slides: the watchdog question, API surface, durable model, example payloads, glossary, 0.x\,$\rightarrow$\,V2 phrasebook.}
    \end{center}
\end{frame}
